% Supplementary Information of the ERL Letter (D-107): copy of methods_dem.tex.
% Methods: waterlines, the static-contour rule, the waterline DEM and emergence dates (A6).
% Design: D-034, D-069, D-075; PREREG section 4 and Amendments 4, 8 and 11.
% Every number comes from \metric{...}.

\subsection*{Waterlines and seabed topography}
The topography follows intertidal elevation models built from satellite waterlines attributed with modelled tide heights\citep{sagar2017,bishoptaylor2019}; here the basin level takes the place of the tide.

\emph{Waterlines.} For each epoch and level bin, the \metric{params.contour_level} contour of water frequency was traced at sub-pixel precision by marching squares\todo{cite: marching squares (Lorensen \& Cline 1987) and scikit-image (van der Walt et al. 2014)}, only through pixels with at least \metric{params.min_valid_obs} valid acquisitions, and sampled at the grid spacing. Each point takes the bin's central level as its height. Its uncertainty combines the level uncertainty of the bin's scenes, without averaging down the part common to scenes gap-filled from Hydroweb in one period, the bin width as a uniform error (width$/\sqrt{12}$) and the residual setup uncertainty (pre-registration section~4 and Amendment~4).

\emph{Static contours.} Where channel banks or ponds rather than the sea set the water's edge, as in the deltas, a contour sits at the same place in every bin and would give one location many heights. For each contour point, $\rho$ is the Spearman correlation at its pixel between water frequency and bin level over the epoch's bins with a frequency there, given at least \metric{params.rho_min_bins} such bins. Sea-controlled shorelines get wetter as the level rises ($\rho$ near 1); banks and ponds do not. The threshold $\rho^*$ is Otsu's split of the histogram of $\rho$ over the E2 and E3 contour points of training blocks, binned as for the scene thresholds. It is used if the split explains at least \metric{params.thr_min_eta} of the variance, has a valley ratio of at most \metric{params.thr_max_valley} and lies between \metric{params.static_rho_lo} and \metric{params.static_rho_hi}; otherwise $\rho^*=\metric{params.static_rho_fallback}$. Points with $\rho<\rho^*$ were excluded from the topography; points without $\rho$ were kept and counted. The rule was fixed on the training blocks of the initial assignment and is not refitted if the blocks grow (pre-registration Amendments~8 and~11). To check that it keeps real steep slopes, where contours of neighbouring bins lie close together, we report the static share of training points within \metric{params.static_steep_m}~m of the next bin's contour beside the overall share. The excluded fraction is reported per sector, split and epoch, and the topography without the exclusion is a robustness check. The observed extents of the hindcast audit are not filtered.

\emph{Gridding.} The retained E2 and E3 points were interpolated linearly on their Delaunay triangulation to a \metric{params.dem_res_m}~m grid in the equal-area projection, inside the stage-2 cover and the triangulation's hull, in \metric{params.dem_tile_km}~km tiles with \metric{params.dem_overlap_km}~km point margins. Every cell stores its interpolation distance, the distance to the nearest contour point. E1 contours serve only the comparison of rising and falling limbs at levels crossed twice. Gaussian-process interpolation is the alternative tested in the robustness checks.

\emph{Emergence dates.} The emergence date of a cell is the first day after the last day on which the daily basin level, interpolated as for the scenes, was at or above the cell's height. It is undefined where the cell was never below the level in the record or has not yet emerged. Because it is read from the cell's own height, a cell placed too low emerges too late.
